\documentclass{article}
\usepackage[margin=1in]{geometry}
\usepackage{amsmath, amssymb, amsthm}
\usepackage{enumitem}
\usepackage{mathtools}
\usepackage{cancel}

% colored links
\usepackage{hyperref}
\hypersetup{
    colorlinks=true,
    linkcolor=blue,
    filecolor=magenta,      
    urlcolor=black!20!BurntOrange,
    }

% Inputting Python code
\usepackage[dvipsnames]{xcolor}
\definecolor{textblue}{rgb}{.2,.2,.7}
\definecolor{textred}{rgb}{0.54,0,0}
\definecolor{textgreen}{rgb}{0,0.43,0}
\usepackage{upquote}
\usepackage{listings}
\lstset{
    language=Python, 
    tabsize=4,
    basicstyle={\ttfamily},
    keywordstyle=\color{textblue},
    commentstyle=\color{textgreen},
    stringstyle=\color{textred},
    frame=none,
    columns=fullflexible,
    keepspaces=true,
    showstringspaces=false,
    xleftmargin=-15mm, % manual adjustment, figure out permanent solution
}

% Drawing figures
\usepackage{tikz}
\usetikzlibrary{calc, shapes.symbols}

% Colored Boxes
\usepackage{tcolorbox}
\tcbuselibrary{skins,hooks}
\usetikzlibrary{shadows}
\usepackage{lipsum}

%Images
\usepackage{graphicx}
    \usepackage{subcaption}
    \usepackage{float}

%Formatting and Spacing
\setitemize[1]{noitemsep, parsep = 5pt, topsep = 5pt}
\setenumerate[1]{label = (\alph*), parsep = 1pt, topsep = 5pt}
\setlength\parindent{0pt}
\linespread{1.1}

%Easy Numbering in case we add or remove questions
\newcounter{counter}
\newcommand{\Exercise}{Exercise \stepcounter{counter}\thecounter}

\newcommand{\Cov}{\text{Cov}}
\newcommand{\trace}{\text{trace}}
\newcommand{\Normal}{\text{Normal}}
\newcommand{\Var}[1]{\operatorname{var}\left(#1\right)}
\newcommand{\E}[1]{\operatorname{E}\left[#1\right]}
% \renewcommand{\Pr}[1]{\operatorname{P}\left(#1\right)}
\newcommand{\oline}[1]{\overleftrightarrow{#1}}

\DeclareMathOperator*{\minimize}{minimize}
\DeclareMathOperator*{\maximize}{maximize}
\DeclareMathOperator*{\argmin}{argmin}
\DeclareMathOperator*{\argmax}{argmax}
\DeclarePairedDelimiter\abs{\lvert}{\rvert}%
\DeclarePairedDelimiter\norm{\lVert}{\rVert}%
\DeclarePairedDelimiterX{\inp}[2]{\langle}{\rangle}{#1, #2}

%Custom Title Fields
\newcommand{\hmwkTitle}{Homework 5 Solutions}
\newcommand{\hmwkDueDate}{Mar.\ 06, 2024}
\newcommand{\hmwkDueTime}{11:59pm EST}
\newcommand{\hmwkClass}{Advanced Algorithms}
\newcommand{\hmwkClassInstructor}{Professor Dana Randall}
\newcommand{\hmwkSection}{Spring 2024}
\newcommand{\hmwkAuthorName}{}

%Headers and Footers
\usepackage{fancyhdr}
\usepackage{extramarks}
\pagestyle{fancy} 
\lhead{CS 4540}
\chead{\hmwkClass \ (\hmwkClassInstructor)}
\rhead{\hmwkTitle}
\cfoot{\thepage}
\renewcommand\headrulewidth{0.4pt}
\renewcommand\footrulewidth{0.4pt}

\title{
    \vspace{2in}
    \textbf{\hmwkClass:\ \hmwkTitle}\\
    \normalsize\vspace{0.1in}\small{Due\ on\ \hmwkDueDate\ at \hmwkDueTime}\\
    \vspace{0.1in}\large{\textit{\hmwkClassInstructor\ \hmwkSection}} \\
    \vspace{0.8in}
    \begin{minipage}[t]{0.4\columnwidth}
        {\footnotesize\textbf{As stated in the syllabus, unauthorized use of previous semester course materials is strictly prohibited in this course.
        }}
    \vspace{0.5in}
    \end{minipage}
    \vspace{0.8in}
    \author{\textbf{\hmwkAuthorName}}
    \date{}
}

\begin{document}

\maketitle
\pagebreak

\subsection*{\Exercise}
  Let us roll a fair 6-sided die $n$ times. Give the best upper bound you can find
  on the probability that the sum of the rolls is at least $5n$.

\subsection*{Solution}
  The mean of a single die roll (1,2,3,4,5,6) is
  $$\E{X}=\sum_{i=1}^6 i \cdot \Pr[X=i] =  \sum_{i=1}^6 i \cdot 1/6 = 7/2.$$ The variance  is $$\Var{X} = \E{X^2} - \E{X}^2 =
  \frac16 (1^2 + 2^2 + 3^2 + 4^2 + 5^2 + 6^2) - (7/2)^2 = 
  \frac{35}{12}.$$ By linearity of expectation, the mean of $n$ die rolls is
  $7n/2$. Because the rolls are independent, the variance is $\frac{35n}{12}$.
  So by Chebychev, $$\Pr(X > 5n) \leq \Pr(\abs{X - 7n/2}) > 3n/2) \leq
  \frac{\frac{35n}{12}}{(3n/2)^2} = \frac{35}{27n}.$$

    \vspace{3mm}
  Markov gives $(7n/2) / 5n = 7/10$. The Markov's inequality bound
  is better when $n=1$, then Chebychev's inequality is better when $n \geq 2$.

    \vspace{3mm}
  The Chernoff bounds given in class are not applicable, since those only
  apply to independent $0,1$ random variables. You would have to prove your
  own version of Chernoff bounds for die rolls.

\subsection*{\Exercise}
    We are trying to predict the outcome of an election. Suppose that there are
  two candidates and the population of N people has exactly N/3 supporters of
  candidate 1 and the rest support the second candidate. If we take a uniform
  sample of n people from this population to poll (with replacement), what is the probability
  that the majority of this n prefer candidate 1? How large should n be to
  guarantee that the probability that candidate 2 wins (in our sample) is at
  least 4/5 according to Markov's inequality? Can you get a better bound on n
  using Chebyshev's inequality? Chernoff bounds?

\subsection*{Solution}
  Let $X$ be the number of people who vote for candidate $1$. Each person
  sampled has a $1/3$ chance of voting for candidate $1$, independently of
  other people sampled, thus this is a binomial distribution, $Bin(n, 1/3)$.
  Immediately, we have: $$\Pr(X > n/2) = \sum_{i > n/2} i \cdot \Pr(X=i) =  \sum_{i > n/2} \binom{n}{i} (1/3)^i (2/3)^{n-i}.$$  However, actually calculating how large $n$ needs to be so that this is at most 1/5 is very challenging, which is where the bounds we've been studying are useful. Note
  $\E{X} = n/3$, and $\Var{X} = 2n/9$, both taken from the binomial
  distribution.

    \vspace{2mm}
  By Markov's inequality, candidate 2 wins in our sample if $X < n/2$: $$\Pr(X
  > n/2) \leq (n/3)/(n/2) = 2/3.$$ This cannot be made lower than 1/5
  
\vspace{2mm}
  By Chebychev's inequality, $$\Pr(X > n/2) \leq \Pr(\abs{X - n/3} > n/6) \leq
  (2n/9) / (n/6)^2 = 8/n.$$ To get this $\leq 1/5$, we need $n \geq 40$
  
\vspace{2mm}
  By Chernoff Bounds, $$\Pr(X > n/2) = \Pr(X > (1 + 1/2)(n/3)) \leq e^{\frac{
  -(1/2)^2 \cdot n/3}{3}} = e^{\frac{n}{36}}.$$
  Using this, we see that $n$ needs to be at least $58$.

 
\subsection*{\Exercise}{
    Consider the following sorting algorithm for $n$ real numbers chosen independently and uniformly from the range $[0,1)$. Place each number $a_i$ in one of the buckets $B_1, B_2, \ldots B_n$ as follows: $a_i$ goes in bucket $B_j$ if $a_i \in \left[\frac{j-1}{n},\frac{j}{n}\right)$. Sort each bucket, (using your favorite sorting algorithm) and output the sorted list.
    \begin{enumerate}
        \item For any integer $0 < M < n$, show that the probability that bucket $B_i$ has at least $M$ entries is at most $1/M$.
        \item Using Chebychev's inequality, give an upper bound on the probability that there exists some bucket with at least $M$ entries.
        \item Can you do better using Chernoff bounds? Give a bound on the probability that there exists some bucket with at least $M$ entries.

        \vspace{3mm}
        \textbf{Note:} Check the restrictions on the Chernoff bound we presented in the notes carefully. Now observe that $\ln(1 + \delta) > \frac{2\delta}{2 + \delta} $ for all \underline{$\delta > 0$}. This implies that $$\delta - (1 + \delta) \ln(1 + \delta) \le \frac{-\delta^{2}}{2 + \delta}.$$
    \end{enumerate}
    
\subsection*{Solution}
    This problem was a bit of a trick question - consequently we were generous
    with the grading.

    \vspace{2mm}
    The over arching idea was: if each bucket has less than a constant $M$ number of elements, then each bucket could be sorted in constant time, and
    we would be left with a $O(n)$ algorithm for sorting We use various probability bounds to calculate the probability that each bucket has less
    than $M$ elements.

    \vspace{2mm}
     The ``trick'' part was this: No matter how you did the bounds, $M$ would have to be a growing function of $n$ to get a probability less than $1$. This was a little vague in the problem statements, so we gave credit based on correctly doing the bounds.

  \begin{enumerate}
    \item A particular bucket has $Bin(n,1/n)$ distribution, which has mean
      $1$. By Markov's inequality, $\Pr(B_j \geq M) \leq 1/M$.
    \item A $Bin(n,1/n)$ distribution has variance $\frac{n-1}{n}$. You also could've calculate this directly using indicator variables (as in problem 2).
      Thus by Chebychev's Inequality, $$\Pr(B_j \geq M) \leq \Pr(\abs{B_j - 1}
      \geq M-1) \leq \frac{n-1}{n(M-1)^2}.$$ Taking a union bound over all
      buckets, the probability that \emph{any} bucket has more than $M$
      elements is $\leq \frac{n-1}{(M-1)^2}$. For this to be $< 1$, $M$ would
      have to be $> \sqrt{n} +1$. (Which means we do not actually sort each
      bucket in constant time.)
    \item In the notes, we obtained the Chernoff bound given in class by first deriving
    $$\begin{aligned}
        \Pr{(X > (1 + \delta) \mu)} &\le e^{(\delta - (1 + \delta)\ln(1 + \delta))\mu}
    \end{aligned}$$
    and then using the fact that
    \begin{equation} \label{eq:1}
        (1 + \delta)\ln(1 + \delta) > \delta + \frac{\delta^{2}}{2} - \frac{\delta^{3}}{6} \ge \delta - \frac{\delta^{2}}{3}
    \end{equation}
    $$.$$ to get $$\Pr{(X > (1 + \delta) \mu)} \le e^{\frac{-\delta^{2}\mu}{3}}\quad (0 < \delta < 1).$$ But, we need a bound that works for all $\delta > 0$. Using the hint, we know that 
    $$(1 + \delta) \ln(1 + \delta) \ge \delta - \frac{-\delta^{2}}{2 + \delta}.$$ Substituting this bound in place of (\ref{eq:1}), we get a alternate version of Chernoff bounds that worked for all $\delta > 0$: $$\Pr(X > (1+\delta)\mu) \leq e^\frac{-\delta^2\mu}{2+\delta}.$$
    We will use this second bound. We first calculate the probability that any given bin has at least $M$ elements in it, where $\mu$ is 1 and $\delta$ is $M-2$: 
    $$\Pr(B_j \geq M) = \Pr( B_j > M-1) = \Pr(B_J > (1 + (M-2))1) \leq e^{-(M-2)^2}{M}.$$ After the union bound, this is
      $ne^\frac{-(M-2)^2}{M}$.  Here, we'd need $M \approx \log(n)$ to get a
      probability $< 1$, which is much better, but still not constant time per
      bucket.
  \end{enumerate}


\subsection*{\Exercise}
 We are given a set of points $\{p_i\}$ on the plane, and we are interested in finding the pair of points that are the farthest apart (called the {\em diameter}). There is an obvious $O(n^2)$ algorithm, but we want an $O(n\log n)$ algorithm.  Prove the following statements and then use them to construct a fast algorithm.

\begin{enumerate}
    \item Prove that the pair of points that are farthest apart are both on the convex hull.

    \vspace{3mm}
    Now assume we are looking for the furthest pair of points on a convex polygon:  $p_1,p_2,\ldots, p_n,$  (given in order going around the polygon).
    \item For two points $p_i,p_j$, we construct the lines $l$ through $p_i$ and $l'$ through $p_j$ so that $l,l'$ are perpendicular to $\overleftrightarrow{p_i,p_j}$. We say that $p_i$ and $p_j$ are \emph{antipodal} if these lines $l,l'$ both do not pass through the convex hull. Show that the pair of points that are the largest distance apart must be antipodal.
    \item Say we are given an edge of the convex hull: $e = (p_1, p_2)$. Give a method to find the point $p_i$ farthest from the line $\overleftrightarrow{p_1,p_2}$. Your method should take $O(i)$ time or faster (not $O(n)$).  (Hint: you'll need the slope of $\overleftrightarrow{p_1,p_2}$, and some basic vector knowledge.)
    \item  Consider adjacent points $p_n,p_1,p_2$, and say that the point farthest from line $\overleftrightarrow{p_n,p_1}$ is $p_i$, and the point farthest from line $\overleftrightarrow{p_1,p_2}$ is $p_j$. Show that $j \geq i$, and the only possible points that are antipodal to $p_1$ are $p_i, p_{i+1}, \ldots p_j$. (Hint: Why are all other points definitely not antipodal?) Now if $p_k$ is farthest from line $\overleftrightarrow{p_2,p_3}$, what about all possible points antipodal to $p_2$?
    \item  Put it all together: Create an algorithm that takes $O(n \log n)$ time to find the pair of pair of points that are farthest apart, with proof of correctness.  Your algorithm should take $O(n)$ steps after finding the convex hull. (Hint: Keep one pointer at $p_1$, and another at $p_i$ as in part d. Alternate updates of $p_i$ and $p_1$ in such a way that all antipodal pairs are considered.  You can't use part c directly, but the idea should help)
\end{enumerate}

\subsection*{Solution}

Note: we will use the phrase convex hull to occasionally mean the boundary
  of the convex hull, especially when we use the preposition ``on'' (i.e.
  ``on'' the convex hull means the boundary, ``in'' the convex hull means the
  interior.)
  \begin{enumerate}[label=(\alph*)]
    \item Let $p_i,p_j$ be the points that are farthest apart. Assume one point, $p_i$, was not on the exterior of the convex hull. Extend the ray $\overrightarrow{p_jp_i}$ until it intersects the boundary of the convex hull at point $x$. If $x$ is in the initial set of points (i.e., is a vertex of the convex hull), then this is a contradiction, as $d(p_i,p_j) < d(x,p_j)$. Alternatively, $x$ is on segment $\overline{p_np_m}$. Then at least one of angles $\angle p_jxp_n$ and $\angle p_jxp_m$ must be at least $90$ degrees. That endpoint is farther from $p_j$ than $x$ (by the Law of Cosines, for instance) and thus farther from $p_j$ than $p_i$, which gives a contradiction in this case too.
    \item Let $p_i,p_j$ be the points that are farthest apart, and label the lines through these points perpendicular to $\overline{p_ip_j}$ as $l_i$ and $l_j$, respectively. Suppose one of these lines, without loss of generality $l_j$, intersects the interior of the convex hull. Then there must be some point $p_k$ on the other side of the line through $l_j$ from $p_i$. Then again, $\angle p_ip_jp_k$ is greater than 90 degrees, and as above, $p_k$ is farther from $p_i$ than $p_j$, which is a contradiction.
            
    \item In constant time, we can find the distance of a point $p_i$ from the line $\oline{p_1p_2}$ in a number of ways. Find this distance for $p_3,p_4,\ldots$ one after another, storing the maximum. If we find point $p_i$ where the distance from $p_{i+1}$ to $\oline{p_1p_2}$ is not larger than the distance from $p_i$ to $\oline{p_1p_2}$, I claim (and prove below) that $p_i$ must be at maximum distance from $\oline{p_1p_2}$, so our algorithm can return $p_i$ without checking the remaining points - it is impossible for the true maximum to be after this local maximum. If $p_i$ is the maximum, we stop after $O(i)$ steps. This is an example of an {\it output-sensitive} running time.

      %We prove this by contradiction: assume that it was possible, and
      %there was some $k > i+1$ such
      %$d(p_k, \oline{p_1p_2}) > 
      %d(p_i, \oline{p_1p_2})$, while at the same time $i$ is a local maximum so that $d(p_i, \oline{p_1p_2}) > d(p_{i+1},
      %\oline{p_1p_2})$.

    \vspace{3mm}
    First, we assume all points on the convex hull are at a unique distance from $\oline{p_1p_2}$ (below we discuss what to do if this is not the case). Orient the polygon so that $\oline{p_1,p_2}$ is horizontal and on the ``bottom'' (all points in the set are above it). This orientation can be used to divide the convex hull into an upper hull and a lower hull (containing $p_1$ and $p_2$). 

    \vspace{3mm}
    Let a point $p_i$ on the convex hull be a {\it local maximum} if both of its neighbors on the convex hull are closer to $\oline{p_1p_2}$ than $p_i$ is. We first show that such a maximum must lie on the upper hull. Suppose, for the sake of contradiction, a local maximum $p_i$ is on the lower hull (and is not one of the endpoints of the lower hull). Then its two neighbors on the
    lower hull are both lower than it (because they are closer to $\oline{p_1p_2}$ by assumption), as is the entire line segment between
    them. This contradicts the fact that $p_i$ is on the convex hull of the points. 

      \vspace{3mm}
      Now, suppose there are two non-consecutive local maxima on the upper hull. Then there must be a point between them on the upper hull that is closer to $\oline{p_1p_2}$ than both. But then this point is below the segment connecting the two maxima, and so it cannot possibly be on the convex hull of the points, a contradiction. 
      Thus, there is at most one local maximum, and so it is a global maximum. 

      \vspace{3mm}
      We haven't yet considered the case where vertices on the convex hull are at the same distance from $\oline{p_1p_2}$. If, during your scan, you encounter a vertex $p_i$ where $p_{i+1}$ is at the same distance from $\oline{p_1p_2}$ as $p_1$, then by the same arguments as above, $p_i$ and $p_{i+1}$ must {\it both} be at maximum distance from $\oline{p_1p_2}$. Either is a satisfactory answer, so your algorithm can just return $p_i$. 
      
     % This only matters in the case when  these vertices are the ones at the maximal distance from the convex hull. In this case, following similar arguments from above, these points must be consecutive along the convex hull, and so the first one you encounter
      
      %but in the case any of those vertices could serve as the global maximum, and we will say by convention that the first one we encounter will. 
      
    %  First, we show by contradiction that any local maximum $p_i$ can't be on the lower hull, where we define a local maximum to be a point on the convex hull whose neighbors on the convex hull are both closer to $\oline{p_1p_2}$. Suppose local maximum $p_i$ is on the lower hull. 
%      Then its two neighbors on the
%      lower hull are both lower than it (because they are closer to $\oline{p_1p_2}$ by assumption), as is the entire line segment between
%      them. This contradicts the fact that $p_i$ is on the lower hull.

    \item 
    %  We state the above result in a slightly different way. The point $p$ of
    %  maximum distance from a line $l$ on a convex body is also the only point
     % such that the line parallel to $l$ through $p$ doesn't intersect the
     % convex body. We use this equivalence below.

    Let $p_i,p_j$, etc. be as in the problem. First, we assume that $p_n$, $p_1$, and $p_2$ are not all collinear, because if they are collinear then we would have $\oline{p_np_1} = \oline{p_1p_2}$ and $j=i$, so we are done. 
    
    \vspace{3mm}
    Assume, for the sake of finding a contradiction, that $j < i$. Orient
    the figure in the direction of $\oline{p_np_1}$ so that
    $p_i$ is the rightmost point, and $\oline{p_np_1}$ is on the left.
    The upper hull of the convex body in this case is
    $p_1 \ldots p_i$, with $p_n$ directly below $p_1$, and by assumption, $p_j$ is on this upper hull. Consider the line $\oline{p_1p_2}$. Note that because there are only right turns along the convex hull, until the total angle of these turns totals 180 degrees each vertex on the upper hull is farther away from $\oline{p_1p_2}$ than the one before it. These angles won't total more than 180 degrees until after you switch from the upper hull to the lower hull at $p_i$, so $p_j$ falls into this realm. This means the vertex after $p_j$ on the upper hull  is farther from $\oline{p_1p_2}$ that $p_j$, a contradiction. Thus, it must be that $j \geq i$. 

    \vspace{3mm}
    All lines through $p_1$ that do not intersect the convex hull must have slope between $\oline{p_np_1}$ and $\oline{p_1p_2}$, otherwise the line would separate those two neighbors. A necessary (but not sufficient) condition for a point $p_h$ to be antipodal to $p_1$ is that there are parallel lines through $p_1$ and $p_h$, neither of which intersect the convex hull.  For any $h < i$, by the arguments above, $p_h$ is closer to both $\oline{p_np_1}$ and $\oline{p_1p_2}$ than $p_i$ is, so any line through $h$ with any slope ranging between the slopes of $\oline{p_np_1}$ and $\oline{p_1p_2}$ (i.e., any line through $p_h$ parallel to a line not intersecting the convex hull through $p_1$) will separate $p_i$ from $p_1$. Thus any such line will intersect the convex hull, meaning $p_1$ and $p_h$ cannot possible be antipodal.  The same argument applies for $h > j$ and the lower hull. Thus the only points that are possibly antipodal to $p_1$ are $p_i, p_{i+1}, \ldots p_j$.

    \vspace{3mm}
    If $p_k$ is the point that is farthest from $\oline{p_2,p_3}$, then all possible points that are antipodal to $p_2$ fall in the range $p_j,..., p_k$. 
    
    \item Now we can state the final algorithm. Find the convex hull, number the points on the hull $p_1$ through $p_n$ in the usual clockwise way.  Keep two pointers, one at $p_1$ and one at $p_i$ as defined above (it takes O(i) time to find i). Increment $i$ until we find $p_j$, the point furthest from $\oline{p_1,p_2}$. Check all the distances $p_1$ with $p_i$ through $p_j$, and keep track of the maximum. Now increment $p_1$ to $p_2$ and repeat, starting with your second pointer already at $p_j$. Both pointers increment
    clockwise around the circle, and at least one of the two pointers increments at each step. Thus we make a complete round of the circle in
    at most $2n$ steps, which takes $O(n)$ time after finding the convex
    hull, for a total $O(n\log n)$ algorithm.
    
    \vspace{3mm}
    Correctness from this algorithm follows from all of the previous parts: the two points that are farthest apart must be antipodal, and the algorithm checks all possible pairs of antipodal points, so it will find the correct solution. 
\end{enumerate}


\end{document}


